\documentclass[11pt]{article}
\usepackage[margin=1in]{geometry}
\usepackage{amsmath,amssymb,booktabs,tabularx,listings,xcolor,hyperref}
\hypersetup{colorlinks=true,linkcolor=blue!50!black}
\lstset{basicstyle=\ttfamily\small,columns=fullflexible,keepspaces=true,breaklines=true,frame=single}
\setlength{\parindent}{0pt}
\setlength{\parskip}{6pt}
\setlength{\emergencystretch}{2em}
\title{CSI 2110 --- Lecture 4\\Trees and Binary Trees}
\author{Study notes based on Bo Sun's Fall 2026 lecture}
\date{}
\begin{document}
\maketitle
\textbf{Source:} \texttt{lec04-tree and binary tree.pdf}, 82 PDF pages. References count PDF pages, not printed slide numbers. Pseudocode includes missing-child and parent-link details; review answers supplement the slides. Heights and depths count \emph{edges}; logarithms are base $2$.

\section{From graphs to rooted trees}
\textit{Source: pp.\ 5--11, 19--20.}

A graph $G=(V,E)$ consists of vertices and edges joining vertex pairs. A path follows adjacent vertices. In the standard undirected-tree setting, a simple path has no repeated vertices, and a cycle closes without repeating intermediate vertices. A graph is connected when every pair of vertices is joined by a path.

A \textbf{tree} is a connected, acyclic graph. It has two key properties:
\begin{enumerate}
\item There is exactly one simple path between any two vertices. Connectivity gives existence; two distinct simple paths would create a cycle.
\item A nonempty tree with $n$ nodes has $n-1$ edges. After choosing a root, each nonroot node contributes exactly one edge to its parent.
\end{enumerate}

Choosing a \textbf{root} establishes a hierarchy:
\begin{itemize}
\item A node's \textbf{parent} is the preceding node on the path from the root; the root has no parent. Its \textbf{children} are the nodes immediately below it.
\item \textbf{Siblings} share a parent. \textbf{Ancestors} are parents, grandparents, and so on; \textbf{descendants} are children, grandchildren, and so on.
\item An \textbf{internal node} has at least one child; an \textbf{external node} or \textbf{leaf} has none.
\item The \textbf{subtree rooted at $v$} contains $v$ and all its descendants.
\item \textbf{Depth} is the number of edges from the root. The root has depth $0$.
\item A node's \textbf{height} is the longest downward path to a leaf. A leaf has height $0$. The tree's height is the root's height, equivalently the maximum node depth.
\end{itemize}

\section{Binary tree families}
\textit{Source: pp.\ 13--17.}

A \textbf{binary tree} has at most two children per node, distinguished as left and right. Even a sole child has a left/right designation.
\begin{center}
\begin{tabularx}{\textwidth}{lX}
\toprule
Type & Additional condition\\
\midrule
Full (proper) & Every node has either zero or two children.\\
Perfect & Full, with all leaves at the same depth.\\
Complete & Every level above the last is full; the last level is filled from left to right without gaps.\\
\bottomrule
\end{tabularx}
\end{center}
A perfect tree is both full and complete. A full tree need not be complete or have logarithmic height. A complete tree can have one node with only a left child, so it need not be full. A binary tree is not automatically a binary \emph{search} tree; ordering of stored keys is a separate property.

Binary decision trees put yes/no questions at internal nodes and outcomes at leaves. Arithmetic-expression trees put binary operators at internal nodes and operands at leaves, naturally giving full binary trees.

\section{Counting properties}
\textit{Source: pp.\ 22--33.}
Let $n$ be the number of nodes, $n_I$ the number of internal nodes, $n_E$ the number of leaves, and $h$ the height. All formulas in this section assume a nonempty tree.

\subsection{Any binary tree}
At depth $i$ there are at most $2^i$ nodes, so
\[
h+1\le n\le 2^{h+1}-1,\qquad
1\le n_E\le2^h,\qquad
h\le n_I\le2^h-1.
\]
A longest root-to-leaf path has $h+1$ nodes, of which $h$ are internal. Summing the maximum level sizes gives the upper bound on $n$. Inverting the bounds gives
\[
\lceil\log_2(n+1)\rceil-1\le h\le n-1.
\]
The ceiling makes the integer lower bound explicit. A chain attains the maximum possible height; a perfect tree attains the maximum node count for a given height.

\subsection{Full binary trees}
Every internal node contributes two child edges. Counting edges both ways gives
\[
2n_I=n-1=n_I+n_E-1\quad\Longrightarrow\quad n_E=n_I+1.
\]
Consequently,
\[
n_E=\frac{n+1}{2},\qquad n_I=\frac{n-1}{2},
\]
and the total node count is odd. Also,
\[
2h+1\le n\le2^{h+1}-1,\qquad h+1\le n_E\le2^h,
\]
\[
\lceil\log_2(n+1)\rceil-1\le h\le\frac{n-1}{2}.
\]
A long chain of internal nodes, each with an extra leaf child, shows why fullness alone does not guarantee a short tree.

\subsection{Perfect and complete trees}
For a perfect tree:
\[
n=2^{h+1}-1,\qquad n_E=2^h,\qquad n_I=2^h-1,
\qquad h=\log_2(n+1)-1.
\]
For a complete tree, the full levels above the last contain $2^h-1$ nodes, and the last level contains between $1$ and $2^h$ nodes:
\[
2^h\le n\le2^{h+1}-1
\quad\Longrightarrow\quad h=\lfloor\log_2 n\rfloor.
\]
\textbf{Example:} at height $3$, a binary tree has $4$--$15$ nodes; a full binary tree has $7$--$15$ nodes with odd total; a complete tree has $8$--$15$ nodes; a perfect tree has exactly $15$.

\section{Tree ADTs}
\textit{Source: pp.\ 35--36.}
The general Tree ADT provides \texttt{root()}, \texttt{parent(p)}, \texttt{children(p)}, \texttt{size()}, \texttt{isEmpty()}, and tests such as \texttt{isRoot}, \texttt{isInternal}, and \texttt{isExternal}. It can expose positions/elements and operations that replace or swap stored elements. \texttt{root()} returns null on an empty tree and \texttt{parent(root)} returns null.

The Binary Tree ADT adds \texttt{leftChild(p)}, \texttt{rightChild(p)}, and \texttt{sibling(p)}. Updates depend on the application. Exchanging stored elements does not exchange the underlying subtrees or change the tree shape.

\section{Traversals}
\textit{Source: pp.\ 37--59, 65--68.}
A traversal visits each node once in a specified order. Fix a left-to-right child order whenever the tree has ordered children.

\subsection{Breadth-first traversal: queue}
\begin{lstlisting}
breadthFirst(root):
    if root = null: return
    Q <- empty queue
    Q.enqueue(root)
    while not Q.isEmpty():
        v <- Q.dequeue()
        visit(v)
        for each child w of v, from left to right:
            Q.enqueue(w)
\end{lstlisting}
The queue processes one level before the next. With constant-time visits and queue operations, time is $\Theta(n)$. Auxiliary space is $O(w)$, where $w$ is maximum level width, hence $O(n)$ in the worst case and $O(1)$ for a chain. The queue can hold parts of two consecutive levels; its size is not always exactly the width of one level.

\subsection{Depth-first traversal: recursion}
For a general rooted tree, preorder visits the root before recursively processing its children; postorder visits it after its children. Binary trees also have inorder:
\begin{center}
\begin{tabular}{ll}
\toprule
Traversal & Visit order\\
\midrule
Preorder & Root, left subtree, right subtree\\
Inorder & Left subtree, root, right subtree\\
Postorder & Left subtree, right subtree, root\\
\bottomrule
\end{tabular}
\end{center}
For a general binary tree, explicitly handle absent children:
\begin{lstlisting}
preOrder(v):
    if v = null: return
    visit(v)
    preOrder(v.left)
    preOrder(v.right)

inOrder(v):
    if v = null: return
    inOrder(v.left)
    visit(v)
    inOrder(v.right)

postOrder(v):
    if v = null: return
    postOrder(v.left)
    postOrder(v.right)
    visit(v)
\end{lstlisting}
All three take $\Theta(n)$ time for constant-time visits. Recursive auxiliary space is $O(h+1)$: $O(\log n)$ for a balanced binary tree and $O(n)$ for a chain. The $n$ total recursive calls do not all coexist on the stack.

\subsection{Worked traversal from the lecture}
The tree on PDF p.\ 46 is:
\begin{lstlisting}
            A
          /   \
         B     G
        / \   / \
       C   D H   I
          / \
         E   F
\end{lstlisting}
\begin{itemize}
\item Breadth-first: A, B, G, C, D, H, I, E, F.
\item Preorder: A, B, C, D, E, F, G, H, I.
\item Inorder: C, B, E, D, F, A, H, G, I.
\item Postorder: C, E, F, D, B, H, I, G, A.
\end{itemize}
For example, inorder completely processes B's subtree before A, so the output begins C, B, E, D, F.

\subsection{Iterative preorder and Euler tour}
An explicit stack implements preorder as follows:
\begin{lstlisting}
S <- empty stack
if root != null: S.push(root)
while not S.isEmpty():
    v <- S.pop()
    visit(v)
    if v.right != null: S.push(v.right)
    if v.left != null: S.push(v.left)
\end{lstlisting}
Push the right child first so that the left child is popped first. For a binary tree the stack requires $O(h+1)$ space.

An \textbf{Euler tour} provides three events at each non-null node: before its left subtree, between its left and right subtrees, and after its right subtree. Recording the first, second, or third event produces preorder, inorder, or postorder respectively.

\section{Applications of traversals}
\textit{Source: pp.\ 42--43, 60--64.}
\begin{itemize}
\item \textbf{Structured document:} preorder prints a heading before its subsections. Depth controls indentation. Traversal overhead is linear; printing indentation also costs time proportional to the amount printed.
\item \textbf{Directory sizes:} postorder obtains child totals before combining them at the parent. Sum file sizes; include directory metadata only if the chosen definition counts it.
\item \textbf{Subtree height:} postorder combines child heights. Using an empty-subtree height of $-1$ handles absent children uniformly.
\end{itemize}
\begin{lstlisting}
height(v):
    if v = null: return -1
    return 1 + max(height(v.left), height(v.right))
\end{lstlisting}
A leaf then has height $1+\max(-1,-1)=0$. Computing the root's height visits all nodes, taking $\Theta(n)$ time.

\subsection{Arithmetic-expression trees}
For $2(a-1)+3b$, preorder gives prefix notation and postorder gives postfix notation:
\begin{lstlisting}
Prefix:   + * 2 - a 1 * 3 b
Postfix:  2 a 1 - * 3 b * +
\end{lstlisting}
Plain inorder does not preserve grouping. Print parentheses around each internal node's expression to obtain $((2\times(a-1))+(3\times b))$.
\begin{lstlisting}
printExpr(v):
    if v is a leaf:
        print(v.element)
    else:
        print("(")
        printExpr(v.left)
        print(v.element)
        printExpr(v.right)
        print(")")

evalExpr(v):
    if v is a leaf: return value of v.element
    x <- evalExpr(v.left)
    y <- evalExpr(v.right)
    return apply(v.element, x, y)
\end{lstlisting}
Evaluation is postorder because operands must be available before applying an operator. For $a=5,b=2$, the result is $2(5-1)+3(2)=14$. Time is $\Theta(n)$ under constant-time arithmetic.

\section{Implementations}
\textit{Source: pp.\ 70--81.}

\subsection{Linked binary tree}
Each node stores its element, parent, left child, and right child. Parent/child access, sibling lookup, and element replacement or swapping take $O(1)$ given valid positions. Sibling lookup must first handle the root, whose parent is null. Storage is $O(n)$.

\texttt{expandExternal(v)} changes a leaf into an internal node with two new children. Set both child links and both children's parent links, and increase size by two. This takes $O(1)$.

The lecture uses \texttt{removeAboveExternal} and later \texttt{removeExternal} for the operation that removes a leaf $v$ \emph{and its parent}, promoting $v$'s sibling $s$:
\begin{enumerate}
\item Require a nonroot leaf whose parent has another child $s$ (as in a full tree).
\item If the parent is the root, set root to $s$ and set \texttt{s.parent = null}.
\item Otherwise, replace the grandparent's link to the parent with $s$, and set \texttt{s.parent = grandparent}.
\item Decrease size by two and invalidate the removed positions.
\end{enumerate}
The operation is $O(1)$. Updating the promoted sibling's parent link is essential even though it is omitted in the nonroot pseudocode on PDF p.\ 74.

\subsection{Array representation}
Number slots from zero:
\[
\operatorname{left}(i)=2i+1,\quad
\operatorname{right}(i)=2i+2,\quad
\operatorname{parent}(i)=\left\lfloor\frac{i-1}{2}\right\rfloor\quad(i>0).
\]
For a complete tree of size $n$, occupied indices are exactly $0,\ldots,n-1$; a child exists only if its index is \emph{less than} $n$. For an arbitrary sparse tree, check array bounds and whether the slot is occupied: the number of nodes is not the array's maximum index.

A right-only chain of $n$ nodes reaches index $2^n-2$, requiring array length $2^n-1$. Thus this representation can waste exponential space on a sparse tree. Promoting a subtree can also require moving all its descendants to different slots. Complete trees, including heaps, avoid these gaps.

\subsection{General trees and first-child/next-sibling encoding}
A general linked tree stores an element, parent, and a collection of child references. Another representation uses two links per node:
\begin{itemize}
\item Left link points to the node's \textbf{first child}.
\item Right link points to its \textbf{next sibling}.
\end{itemize}
The root is unchanged. Traversing right links enumerates siblings; it does not descend to children of the original general tree. This encoding preserves preorder with consistent child ordering, but original depths and parent relationships cannot be read as ordinary binary-tree depths and parents.

\section{Quick review}
\begin{enumerate}
\item A full tree has $10$ internal nodes. How many leaves and total nodes?\\
\textbf{Answer:} $11$ leaves and $21$ nodes.
\item A complete tree has $12$ nodes. What is its height?\\
\textbf{Answer:} $\lfloor\log_2 12\rfloor=3$.
\item Which traversal evaluates an expression tree?\\
\textbf{Answer:} postorder, because child values are computed before the operator.
\item Why must iterative preorder push right before left?\\
\textbf{Answer:} the stack is LIFO, so left must be pushed last to be visited next.
\item Does inorder always output sorted keys?\\
\textbf{Answer:} no. That requires the binary-search-tree ordering property, which an arbitrary binary tree does not have.
\end{enumerate}
\end{document}
